\documentclass[12pt, ukrainian]{article}
\usepackage[a4paper]{geometry}\geometry{top=1.5cm,bottom=1.5cm,left=2.5cm,right=1.5cm}%\setlength\textwidth{190mm} \setlength\textheight{237mm}%\setlength\voffset{-18mm} \setlength\hoffset{-33mm}\usepackage[T2A]{fontenc}
\usepackage[utf8]{inputenc}
\usepackage[ukrainian]{babel}
\begin{document}
{\Large \center  Варіанти завдань до програмного проєкту №3, СА\par 2023/24 н.р.\par \bigskip\bigskip}

«Зростаюча серія» – це послідовність чисел, що йдуть поспіль і неспадають, наприклад 2 2 2 5 7 7.

«Власна зростаюча серія» – це зростаюча серія, що зростає, наприклад 2  5  7.

«Спадна серія» – це послідовність чисел, що йдуть поспіль і незростають, наприклад 7 7 5 3 3 2.

«Власна спадна серія» – це спадна серія, що спадає, наприклад~7~5~3.

«Серія» – це зростаюча або спадна серія.

«Власна серія» – це власна зростаюча або власна спадна серія.

«Пилка» – це послідовність чисел, що йдуть поспіль і чергуються за неспаданням та незростанням, наприклад 2 2 2 5 3 7 4 6 5 9 (завершуватися може і неспаданням, і незростанням).

«Парна пилка» – це пилка, що завершується незростанням, наприклад 2 2 2 5 3.

«Непарна пилка» – це пилка, що завершується неспаданням, наприклад 2 2 2 5 3 7.

«Власна пилка» – це послідовність чисел, що йдуть поспіль і чергуються за зростанням та спаданням, наприклад 2 5 3 7 4 6 5 9 (завершуватися може і зростанням, і спаданням).

«Парна власна пилка» – це власна пилка, що завершується спаданням, наприклад 2 5 3 7 6.

«Непарна власна пилка» – це власна пилка, що завершується зростанням, наприклад 2 5 3 7.

\textbf{Увага!} 1. У пилці після неспадання обов’язково має бути незростання, а після незростання  – неспадання. Тобто послідовності 2 5 7 та 2 7 5 3 не утворюють пилку. Аналогічне стосується й власної пилки. 2. Пилка має починатися неспаданням, власна пилка – зростанням.

«Пік» – це послідовність чисел, що йдуть поспіль і спочатку зростають, а потім спадають, наприклад 2 3 5 4 3. 

«Воронка» – це послідовність чисел, що йдуть поспіль і спочатку спадають, а потім зростають, наприклад 5 3 2 3 4.

«Плато» – це послідовність чисел, що йдуть поспіль і спочатку зростають, потім повторюються, а потім спадають, наприклад 2 3 5 5 4 3.

«Яма» – це послідовність чисел, що йдуть поспіль і спочатку спадають, потім повторюються, а потім зростають, наприклад 5 3 2 2 3 4.

«Гора» – це послідовність чисел, що йдуть поспіль і спочатку не спадають, а потім не зростають, наприклад 2 2 3 3 5 7 7 6 6 4 3 3. «Гора» починається неспаданням, а завершується незростанням. 

«Власна передгора» – це гора, що починається зростанням, наприклад 2 3 3 5 7 6 6 4 3 3. 

«Власна післягора» – це гора, що завершується спаданням, наприклад  3 5 7 7 6 6 4 3.

«Власна гора» – це гора, що починається зростанням, а завершується неспаданням.

«Долина» – послідовність чисел, що йдуть поспіль і спочатку не зростають, а потім не спадають, наприклад 7 7 6 6 3 3 4 5 5 6 6. «Долина» починається незростанням, а завершується неспаданням. 

«Власна переддолина» – це долина, що починається спаданням, наприклад 7 6 6 3 3 4 5 5 6 6. 

«Власна післядолина» – це долина, що завершується зростанням, наприклад 6 6 3 3 4 5 5 6.

\par \bigskip
\newpage

{\center  Варіанти \par \bigskip\bigskip}

%begin
101) \space 
З послідовності цілих вилучили всі числа, у трійковому запису яких присутня цифра~2. Знайти та вивести першу найдовшу підпослідовність залишку, що є власною переддолиною та має довжину не менше двох.
%власною переддолиною
%у трійковому запису яких присутня цифра~2

\par
%end
\vspace{0.9cm}
%begin
102) \space 
З послідовності цілих вилучили всі числа, у сімковому запису яких присутня цифра~3. Знайти та вивести першу найдовшу підпослідовність залишку, що є горою та має довжину не менше двох.
%горою
%у сімковому запису яких присутня цифра~3

\par
%end
\vspace{0.9cm}
%begin
103) \space 
З послідовності цілих вилучили всі числа, у сімковому запису яких відсутня цифра~3. Знайти та вивести першу найдовшу підпослідовність залишку, що є серією та має довжину не менше двох.
%серією
%у сімковому запису яких відсутня цифра~3

\par
%end
\vspace{0.9cm}
%begin
104) \space 
З послідовності цілих вилучили всі числа, які мають більше чотирьох різних простих дільників. Знайти та вивести першу найдовшу підпослідовність залишку, що є власною післядолиною та має довжину не менше двох.
%власною післядолиною
%які мають більше чотирьох різних простих дільників

\par
%end
\vspace{0.9cm}
%begin
105) \space 
З послідовності цілих вилучили всі числа, у трійковому запису яких нема парних цифр. Знайти та вивести першу найдовшу підпослідовність залишку, що є піком та має довжину не менше двох.
%піком
%у трійковому запису яких нема парних цифр

\par
%end
\vspace{0.9cm}
%begin
106) \space 
З послідовності цілих вилучили всі числа, у трійковому запису яких нема цифри~1. Знайти та вивести першу найдовшу підпослідовність залишку, що є плато та має довжину не менше двох.
%плато
%у трійковому запису яких нема цифри~1

\par
%end
\vspace{0.9cm}
%begin
107) \space 
З послідовності цілих вилучили всі числа, у сімковому запису яких присутня цифра~5. Знайти та вивести першу найдовшу підпослідовність залишку, що є парною власною пилкою та має довжину не менше двох.
%парною власною пилкою
%у сімковому запису яких присутня цифра~5

\par
%end
\vspace{0.9cm}
%begin
108) \space 
З послідовності цілих вилучили всі числа, які мають не більше п'яти різних простих дільників. Знайти та вивести першу найдовшу підпослідовність залишку, що є долиною та має довжину не менше двох.
%долиною
%які мають не більше п'яти різних простих дільників

\par
%end
\vspace{0.9cm}
%begin
109) \space 
З послідовності цілих вилучили всі числа, у трійковому запису яких нема непарних цифр. Знайти та вивести першу найдовшу підпослідовність залишку, що є власною спадною серією та має довжину не менше двох.
%власною спадною серією
%у трійковому запису яких нема непарних цифр

\par
%end
\vspace{0.9cm}
%begin
110) \space 
З послідовності цілих вилучили всі числа, у сімковому запису яких нема парних цифр. Знайти та вивести першу найдовшу підпослідовність залишку, що є непарною власною пилкою та має довжину не менше двох.
%непарною власною пилкою
%у сімковому запису яких нема парних цифр

\par
%end
\newpage
%begin
111) \space 
З послідовності цілих вилучили всі числа, у сімковому запису яких нема непарних цифр. Знайти та вивести першу найдовшу підпослідовність залишку, що є парною пилкою та має довжину не менше двох.
%парною пилкою
%у сімковому запису яких нема непарних цифр

\par
%end
\vspace{0.9cm}
%begin
112) \space 
З послідовності цілих вилучили всі числа, які мають більше п'яти різних простих дільників. Знайти та вивести першу найдовшу підпослідовність залишку, що є зростаючою серією та має довжину не менше двох.
%зростаючою серією
%які мають більше п'яти різних простих дільників

\par
%end
\vspace{0.9cm}
%begin
113) \space 
З послідовності цілих вилучили всі числа, у трійковому запису яких присутня цифра~1. Знайти та вивести першу найдовшу підпослідовність залишку, що є власною горою та має довжину не менше двох.
%власною горою
%у трійковому запису яких присутня цифра~1

\par
%end
\vspace{0.9cm}
%begin
114) \space 
З послідовності цілих вилучили всі числа, які мають не більше чотирьох різних простих дільників. Знайти та вивести першу найдовшу підпослідовність залишку, що є власною післягорою та має довжину не менше двох.
%власною післягорою
%які мають не більше чотирьох різних простих дільників

\par
%end
\vspace{0.9cm}
%begin
115) \space 
З послідовності цілих вилучили всі числа, у трійковому запису яких нема цифри~2. Знайти та вивести першу найдовшу підпослідовність залишку, що є власною серією та має довжину не менше двох.
%власною серією
%у трійковому запису яких нема цифри~2

\par
%end
\vspace{0.9cm}
%begin
116) \space 
З послідовності цілих вилучили всі числа, у сімковому запису яких відсутня цифра~5. Знайти та вивести першу найдовшу підпослідовність залишку, що є спадною серією та має довжину не менше двох.
%спадною серією
%у сімковому запису яких відсутня цифра~5

\par
%end
\vspace{0.9cm}
%begin
117) \space 
З послідовності цілих вилучили всі числа, у сімковому запису яких нема цифри~2. Знайти та вивести першу найдовшу підпослідовність залишку, що є власною зростаючою серією та має довжину не менше двох.
%власною зростаючою серією
%у сімковому запису яких нема цифри~2

\par
%end
\vspace{0.9cm}
%begin
118) \space 
З послідовності цілих вилучили всі числа, у трійковому запису яких нема парних цифр. Знайти та вивести першу найдовшу підпослідовність залишку, що є непарною пилкою та має довжину не менше двох.
%непарною пилкою
%у трійковому запису яких нема парних цифр

\par
%end
\vspace{0.9cm}
%begin
119) \space 
З послідовності цілих вилучили всі числа, які мають не більше п'яти різних простих дільників. Знайти та вивести першу найдовшу підпослідовність залишку, що є власною пилкою та має довжину не менше двох.
%власною пилкою
%які мають не більше п'яти різних простих дільників

\par
%end
\vspace{0.9cm}
%begin
120) \space 
З послідовності цілих вилучили всі числа, у сімковому запису яких відсутня цифра~3. Знайти та вивести першу найдовшу підпослідовність залишку, що є власною передгорою та має довжину не менше двох.
%власною передгорою
%у сімковому запису яких відсутня цифра~3

\par
%end
\newpage
%begin
121) \space 
З послідовності цілих вилучили всі числа, які мають більше п'яти різних простих дільників. Знайти та вивести першу найдовшу підпослідовність залишку, що є воронкою та має довжину не менше двох.
%воронкою
%які мають більше п'яти різних простих дільників

\par
%end
\vspace{0.9cm}
%begin
122) \space 
З послідовності цілих вилучили всі числа, у трійковому запису яких нема цифри~1. Знайти та вивести першу найдовшу підпослідовність залишку, що є ямою та має довжину не менше двох.
%ямою
%у трійковому запису яких нема цифри~1

\par
%end
\vspace{0.9cm}
%begin
123) \space 
З послідовності цілих вилучили всі числа, у сімковому запису яких нема парних цифр. Знайти та вивести першу найдовшу підпослідовність залишку, що є власною зростаючою серією та має довжину не менше двох.
%власною зростаючою серією
%у сімковому запису яких нема парних цифр

\par
%end
\vspace{0.9cm}
%begin
124) \space 
З послідовності цілих вилучили всі числа, які мають не більше чотирьох різних простих дільників. Знайти та вивести першу найдовшу підпослідовність залишку, що є власною передгорою та має довжину не менше двох.
%власною передгорою
%які мають не більше чотирьох різних простих дільників

\par
%end
\vspace{0.9cm}
%begin
125) \space 
З послідовності цілих вилучили всі числа, у сімковому запису яких відсутня цифра~5. Знайти та вивести першу найдовшу підпослідовність залишку, що є піком та має довжину не менше двох.
%піком
%у сімковому запису яких відсутня цифра~5

\par
%end
\vspace{0.9cm}
%begin
126) \space 
З послідовності цілих вилучили всі числа, які мають більше чотирьох різних простих дільників. Знайти та вивести першу найдовшу підпослідовність залишку, що є плато та має довжину не менше двох.
%плато
%які мають більше чотирьох різних простих дільників

\par
%end
\vspace{0.9cm}
%begin
127) \space 
З послідовності цілих вилучили всі числа, у сімковому запису яких присутня цифра~5. Знайти та вивести першу найдовшу підпослідовність залишку, що є парною власною пилкою та має довжину не менше двох.
%парною власною пилкою
%у сімковому запису яких присутня цифра~5

\par
%end
\vspace{0.9cm}
%begin
128) \space 
З послідовності цілих вилучили всі числа, у сімковому запису яких присутня цифра~3. Знайти та вивести першу найдовшу підпослідовність залишку, що є парною пилкою та має довжину не менше двох.
%парною пилкою
%у сімковому запису яких присутня цифра~3

\par
%end
\vspace{0.9cm}
%begin
129) \space 
З послідовності цілих вилучили всі числа, у трійковому запису яких нема непарних цифр. Знайти та вивести першу найдовшу підпослідовність залишку, що є долиною та має довжину не менше двох.
%долиною
%у трійковому запису яких нема непарних цифр

\par
%end
\vspace{0.9cm}
%begin
130) \space 
З послідовності цілих вилучили всі числа, у трійковому запису яких нема цифри~2. Знайти та вивести першу найдовшу підпослідовність залишку, що є власною післядолиною та має довжину не менше двох.
%власною післядолиною
%у трійковому запису яких нема цифри~2

\par
%end
\newpage
%begin
131) \space 
З послідовності цілих вилучили всі числа, у трійковому запису яких присутня цифра~2. Знайти та вивести першу найдовшу підпослідовність залишку, що є власною пилкою та має довжину не менше двох.
%власною пилкою
%у трійковому запису яких присутня цифра~2

\par
%end
\vspace{0.9cm}
%begin
132) \space 
З послідовності цілих вилучили всі числа, у сімковому запису яких нема цифри~2. Знайти та вивести першу найдовшу підпослідовність залишку, що є горою та має довжину не менше двох.
%горою
%у сімковому запису яких нема цифри~2

\par
%end
\vspace{0.9cm}
%begin
133) \space 
З послідовності цілих вилучили всі числа, у сімковому запису яких нема непарних цифр. Знайти та вивести першу найдовшу підпослідовність залишку, що є спадною серією та має довжину не менше двох.
%спадною серією
%у сімковому запису яких нема непарних цифр

\par
%end
\vspace{0.9cm}
%begin
134) \space 
З послідовності цілих вилучили всі числа, у трійковому запису яких присутня цифра~1. Знайти та вивести першу найдовшу підпослідовність залишку, що є власною серією та має довжину не менше двох.
%власною серією
%у трійковому запису яких присутня цифра~1

\par
%end
\vspace{0.9cm}
%begin
135) \space 
З послідовності цілих вилучили всі числа, які мають не більше п'яти різних простих дільників. Знайти та вивести першу найдовшу підпослідовність залишку, що є непарною власною пилкою та має довжину не менше двох.
%непарною власною пилкою
%які мають не більше п'яти різних простих дільників

\par
%end
\vspace{0.9cm}
%begin
136) \space 
З послідовності цілих вилучили всі числа, у трійковому запису яких нема цифри~1. Знайти та вивести першу найдовшу підпослідовність залишку, що є зростаючою серією та має довжину не менше двох.
%зростаючою серією
%у трійковому запису яких нема цифри~1

\par
%end
\vspace{0.9cm}
%begin
137) \space 
З послідовності цілих вилучили всі числа, у трійковому запису яких нема непарних цифр. Знайти та вивести першу найдовшу підпослідовність залишку, що є власною спадною серією та має довжину не менше двох.
%власною спадною серією
%у трійковому запису яких нема непарних цифр

\par
%end
\vspace{0.9cm}
%begin
138) \space 
З послідовності цілих вилучили всі числа, у сімковому запису яких нема парних цифр. Знайти та вивести першу найдовшу підпослідовність залишку, що є ямою та має довжину не менше двох.
%ямою
%у сімковому запису яких нема парних цифр

\par
%end
\vspace{0.9cm}
%begin
139) \space 
З послідовності цілих вилучили всі числа, у трійковому запису яких нема парних цифр. Знайти та вивести першу найдовшу підпослідовність залишку, що є непарною пилкою та має довжину не менше двох.
%непарною пилкою
%у трійковому запису яких нема парних цифр

\par
%end
\vspace{0.9cm}
%begin
140) \space 
З послідовності цілих вилучили всі числа, у сімковому запису яких присутня цифра~3. Знайти та вивести першу найдовшу підпослідовність залишку, що є воронкою та має довжину не менше двох.
%воронкою
%у сімковому запису яких присутня цифра~3

\par
%end
\newpage
\end{document}
